\chapter{Methods}\label{ch_methods}

The preceding chapters established that no existing tool jointly optimises map styles and tile data, and introduced the compiler, database, and encoding techniques whose analogues make such joint optimisation possible.
This chapter applies those foundations to the design and implementation of a concrete optimiser: a standalone Rust tool that consumes a MapLibre style~JSON document and, optionally, tile-level statistics, and produces a semantically equivalent style together with optimised tile data.
All transformations are \emph{visually lossless}: the optimised output must produce pixel-identical rendering for every tile at every zoom level, satisfying the visual equivalence constraint formalised in \cref{sec_problem_statement}.

\section{Scope}\label{sec_scope}

As discussed in \cref{sec_styles}, the most common target for performance research on interactive maps is a system built on a rich, dynamic style specification.
Two implementations dominate this space: Mapbox~styles and MapLibre~styles.
Since the Mapbox codebase is proprietary and evaluations may not be performed without explicit consent, only MapLibre is a viable research target.

Among MapLibre renderers, the web implementation (\texttt{maplibre-gl-js}) offers the most complete support for the styling language and is the most widely used \emph{as an interactive developer target}: the mobile MapLibre Native renderers share much of the styling engine, but the web renderer is the reference implementation that first tracks new specification features, exposes the richest instrumentation hooks for Puppeteer-driven benchmarking, and runs on the same hardware profile as the benchmark host.
It was therefore selected as the system under test for all rendering benchmarks.

The optimiser itself is a standalone Rust binary.
It takes a style~JSON file as input, applies a configurable set of optimisations passes, and writes an optimised style~JSON as output.
When tile statistics are provided (via an MBTiles database or a pre-computed statistics file), additional data-driven passes become available.
No integration with a specific tile server is required; the optimised artefacts can be served by any standards-compliant tile infrastructure.
For more realistic results, a network delay is implemented to avoid artificially high numbers due to local network transfer.

The tile re-encoder described in \cref{sub_data_optimisations} lives in the same Cargo workspace as the style optimiser, but is deliberately split into three crates so that it can also be consumed from a \ac{JVM}-based tile-generation pipeline without re-implementing the encoder.
The core encoder and decoder live in \texttt{rust/mlt-core}; a thin C~\ac{ABI} shim in \texttt{rust/mlt-ffi} exposes a stable surface whose header is generated via \texttt{cbindgen} and checked in as \texttt{mlt\_ffi.h}; and \texttt{rust/mlt-java} wraps that shim as a Java library using the Panama \ac{FFM} \ac{API} (JEP~454), with Java bindings regenerated from the header via \texttt{jextract}.
Panama was chosen over conventional \ac{JNI} glue because it removes the need for hand-written C stubs and gives the \ac{JVM} direct, zero-copy access to off-heap buffers - exactly the boundary shape that a columnar encoder needs, because each tile is handed over as a batch of parallel \texttt{MemorySegment}s rather than one call per feature, amortising the per-boundary cost over thousands of features in one encoder invocation.
The resulting artefact is published to Maven Central as \texttt{org.maplibre:mlt-encoder-native}, shipped either as a single fat jar that bundles native libraries for all supported targets or as per-platform slim jars (\texttt{linux-x86\_64}, \texttt{linux-x86\_64-musl}, \texttt{linux-aarch64}, \texttt{macos-aarch64}, \texttt{windows-x86\_64}), and requires Java~25 or later for stable Panama support.
This wrapper is the intended Planetiler integration path: Planetiler already builds tiles in a \ac{JVM} process, so routing its in-memory \ac{JTS} geometries and property tables through \texttt{MltEncoder.encode(\ldots)} is cheaper than serialising them to \ac{MVT} protobuf and re-parsing in a separate Rust process.
The legacy self-contained Java encoder in \texttt{java/mlt-core}~\cite{tremmelMapLibreTileNext2025} is kept only as a reference baseline for the \ac{MLT} experiments in \cref{ch_experiments} and is not a target of the optimisations proposed in this thesis.

\section{Approach}\label{sec_approach}

To evaluate the effectiveness of the optimisations strategies developed in this work, a benchmarking-driven approach was adopted.
A reliable performance baseline was established first, followed by cumulative ablation experiments that isolate the contribution of each optimisations pass.

\subsection{Two-Phase Optimisation Pipeline}\label{sub_pipeline}

The optimiser operates in two distinct phases, each targeting a different level of abstraction:

\begin{enumerate}
    \item \textbf{Expression passes} operate directly on the JSON representation of style expressions.
    They apply peephole rewrites, constant folding, and algebraic simplifications to individual expressions without deserialising the style into typed Rust structures.
    Expression passes run to a fixpoint: the pass set is applied repeatedly (up to eight iterations) until no further changes are detected.

    \item \textbf{Structural passes} deserialise the style into typed Rust structures generated from the upstream MapLibre specification, apply transformations that require cross-layer or cross-source reasoning, and serialise the result back to JSON.
    These passes handle dead layer elimination, metadata refinement, layer merging, and cleanup.
\end{enumerate}

After the structural phase, a post-processing stage prunes zoom-bounded interpolation and step stops that have become unreachable due to tightened metadata bounds.
A final fixpoint re-fold cleans up residual trivial expressions (e.g., \texttt{["all",~x]} $\to$ \texttt{x}) introduced by predicate removal.
Layer merging runs last, after dead layers have been eliminated, to maximise the number of mergeable adjacent layers.

The separation into two phases is deliberate.
Expression passes operate on raw JSON and require no knowledge of the style schema beyond expression syntax.
This permits fine-grained, composable rules that each match a local pattern and apply a rewrite; the rules need not be aware of one another, and new rules can be added without modifying existing ones.
Structural passes, by contrast, require typed access to the full style document - for example, dead layer elimination must resolve source-layer references across layers and sources, and layer merging must compare paint properties across layers.
Deserialising the style into typed Rust structures incurs a non-trivial cost (parsing, validation, allocation), so it is performed exactly once, after expression passes have simplified the expressions that the typed structures will contain.

Fixpoint iteration in the expression phase is necessary because rewrites compose non-trivially: constant folding a \texttt{has} expression may produce a boolean literal that enables dead-branch elimination in an enclosing \texttt{case}, which may in turn expose a boolean absorption opportunity in the filter's top-level \texttt{all} node.
The cap of eight iterations provides a defensive safety margin; in practice, most styles converge within two to three iterations, and no style in either the deep benchmarking corpus or the Maputnik generalisation corpus has ever reached the cap.

A notable cross-phase optimisations is \emph{filter-to-property constant propagation}.
During the typed structural phase, equality constraints are extracted from a layer's filter expression (e.g.\ if the filter asserts \texttt{["==",~["get",~"class"],~"primary"]}, then the property \texttt{class} is known to equal \texttt{"primary"} for all features passing the filter).
These constraints are propagated into the layer's paint and layout expressions, replacing \texttt{["get",~"class"]} references with the known literal value.
This enables further constant folding within paint ramp expressions that branch on the same property, collapsing \texttt{match} and \texttt{case} arms that are unreachable given the filter constraint.
The propagation operates within a single layer (intra-procedural, in compiler terminology) and extracts constraints from \texttt{all}, equality, \texttt{has}, \texttt{in}, and range predicates.

\subsection{Code Generation from v8.json}\label{sub_codegen}

Type-safe manipulation of style documents is central to the correctness guarantees of the structural passes.
Rather than operating on raw JSON, the optimiser deserialises styles into Rust structures that are \emph{generated} from the upstream MapLibre GL~JS style specification.
The generation pipeline proceeds in three stages.
First, the machine-readable style specification published by the MapLibre project is \emph{decoded} into an internal representation.
Second, the decoded data is \emph{normalised} into a stable intermediate representation that resolves cross-cutting concerns -- layers combine data from multiple specification sections, expressions require special handling, and so forth.
Third, the intermediate representation is consumed by a code emitter that \emph{generates} Rust source files with serialisation and deserialisation support.

This pipeline ensures that every structural pass operates on types that faithfully mirror the specification, and that changes to the upstream spec can be incorporated by re-running the generator.
The alternative - hand-written Rust types - would require manual synchronisation with specification changes and risk silent divergence.
The generated types include serialisation and deserialisation support, so the structural passes can round-trip through JSON without losing unknown fields that the generator has not yet modelled.

The normalisation stage produces what we call the \texttt{MIR} (mid-level intermediate representation), which resolves cross-cutting concerns in the specification.
For example, the upstream \texttt{v8.json} defines layer properties across multiple sections (paint, layout) and references expression operators in a separate section; the MIR collapses these into per-layer field definitions that include the property's type, default value, and expression capabilities (whether the property supports zoom-dependent, feature-dependent, or both kinds of expressions).
This metadata is critical for structural passes: layer merging queries the MIR to determine whether a differing paint property supports feature-driven expressions (and can therefore be wrapped in a \texttt{case}/\texttt{match} over the filter discriminator) or is camera-only (and must be uniform across merged layers).
The MIR also drives the default stripping pass, which compares property values against the specification defaults encoded in the generated types.